\documentclass[journal=cmatex,manuscript=article,layout=twocolumn]{achemso}

\usepackage{graphicx}
\usepackage{booktabs}
\usepackage{amsmath}

\graphicspath{%
  {D:/DINOSR/Claude/PaperRun_claude/dino_sr_contrastive/runs/_sweep_m_K_20260525_213539/IMC_SI5/stage2/m0.9700_seed42_K60/_interpretability/}%
  {D:/DINOSR/Claude/PaperRun_claude/dino_sr_contrastive/runs/_sweep_m_K_20260525_213539/IMC_SI5/stage2/m0.9700_seed42_K60/eval/paper_attribution/}%
}

\author{Daniel Kahn}
\email{author@institution.edu}
\affiliation{Department of Materials Science and Engineering, [Institution]}

\title{Physical Basis of Self-Supervised Clustering of 4D-STEM
Diffraction: Scattered Intensity and Low-$q$ Structure, Not
Crystallographic Orientation}

\abbreviations{4D-STEM, NBED, ACOM, DINO, SupCon}
\keywords{4D-STEM, nanobeam electron diffraction, self-supervised learning,
interpretability, automated crystal orientation mapping}

\begin{document}

\begin{abstract}
Self-supervised clustering of four-dimensional scanning transmission
electron microscopy (4D-STEM) diffraction data groups scan positions into
classes without labels, but the \emph{physical} meaning of those classes is
typically opaque. We interrogate a validated DINO model---trained with a
radial-gated supervised-contrastive objective and azimuthal ($\theta$-roll)
augmentation---on nanobeam electron diffraction (NBED) of a 150~nm
indomethacin (IMC) molecular-crystal film. Combining linear probes,
correlation-ratio analysis, causal input ablations with re-inference,
automated crystal orientation mapping (ACOM), and gradient attribution, we
find the $K=13$ partition is governed by two physical axes: (i) the total
scattered (diffracted) intensity---a coupled crystallinity/mass-thickness
signal that dominates the coarse grouping---and (ii) the two-dimensional
low-$q$ diffraction structure that resolves the fine classes. The partition
is invariant to in-plane crystallographic orientation (DINO vs.\ ACOM zone
axis, adjusted mutual information $\mathrm{AMI}=0.05$) and independent of the
raw Bragg-peak count, consistent with the rotation-invariant training design.
\end{abstract}

\section{Introduction}

Unsupervised and self-supervised analysis of 4D-STEM data is attractive: it
can segment a scan into structurally distinct regions without a physical
forward model or human labels.\cite{caron2021,savitzky2021} The cost is
interpretability---a learned class map is only scientifically actionable once
we know \emph{which} physical property separates the classes. Here we treat
one validated model as the object of study and ask, quantitatively, what it
keys on. The model is a self-distillation (DINO) encoder\cite{caron2021}
augmented with a radial-gated supervised-contrastive loss\cite{khosla2020}
and azimuthal ($\theta$-roll) augmentation, the latter enforcing invariance
to in-plane rotation of the diffraction pattern. We apply a four-part
interpretability battery---representational probing, causal ablation,
orientation cross-correlation, and spatial attribution---to a 150~nm
indomethacin film imaged in NBED mode.

\section{Methods}

\textbf{Data and model.} A $128\times128$ scan of NBED patterns
($192\times192$ detector, intensity clipped at $v_{\max}=5$) was clustered by
the trained model into $K_{\mathrm{active}}=13$ classes. All analyses below
reuse the frozen checkpoint and its saved inference (per-position embeddings,
soft assignments, and hard labels).

\textbf{(1) Representational probing.} For each candidate physical factor we
report a 5-fold cross-validated ridge-regression $R^2$ (\emph{how linearly
decodable} the factor is from the embedding) and the correlation ratio
$\eta^2$ (\emph{how much} of the factor's variance the DINO partition
explains). Factors: total scattered intensity (post-beam energy),
crystallinity (1-D peak/halo ratio), spottiness (azimuthal variance), ACOM
correlation, and Bragg-peak count; categorical ACOM phase is scored by AMI
and the adjusted Rand index (ARI).

\textbf{(2) Causal ablation.} We perturb the \emph{input} (removing one
physical degree of freedom at a time), re-infer with the same model, and
measure ARI against the original map. High ARI $\Rightarrow$ that information
was not needed; low ARI $\Rightarrow$ the model depends on it.

\textbf{(3) Orientation cross-check.} We index every $4^{\mathrm{th}}$ scan
position by ACOM (py4DSTEM\cite{savitzky2021,ophus2022}) using an
\texttt{auto} (symmetry-derived) orientation plan over the $\alpha$ and
$\gamma$ IMC polymorphs ($k_{\max}=0.5$~\AA$^{-1}$, peak detection threshold
$0.1$, GPU-accelerated), then correlate DINO classes against the recovered
zone axis.

\textbf{(4) Spatial attribution.} Gradient-weighted class activation mapping
(Grad-CAM) and integrated gradients (IG) on the per-class average pattern,
projected from the polar input back onto the Cartesian detector.

\section{Results and Discussion}

\subsection{What the embedding encodes}

Table~\ref{tab:probe} and Figure~\ref{fig:rank} rank the factors. The
embedding most strongly represents the \emph{total scattered intensity}
(probe $R^2=0.71$, $\eta^2=0.34$), followed by \emph{crystallinity}
($R^2=0.24$). Crystallographic observables---ACOM correlation and Bragg-peak
count---are essentially undecodable ($R^2\approx0$, $\eta^2<0.01$), and the
categorical ACOM phase is uncorrelated with the partition
($\mathrm{AMI}\approx0$). This is correlational; the ablations below test
causality.

\begin{table}[t]
\centering
\caption{Representational probing. Probe $R^2$: 5-fold CV ridge regression
embedding $\rightarrow$ factor. $\eta^2$: correlation ratio of the DINO
partition.}
\label{tab:probe}
\begin{tabular}{lcc}
\toprule
Factor & Probe $R^2$ & $\eta^2$ \\
\midrule
Scattered intensity        & \textbf{0.71} & \textbf{0.34} \\
Crystallinity (peak/halo)  & \textbf{0.24} & 0.15 \\
Spottiness (azim.\ var.)   & 0.15          & 0.02 \\
ACOM correlation           & $\approx 0$   & 0.01 \\
Bragg-peak count           & $\approx 0$   & 0.01 \\
ACOM phase ($\alpha/\gamma$) & ---         & AMI $\approx 0$ \\
\bottomrule
\end{tabular}
\end{table}

\begin{figure}[t]
\centering
\includegraphics[width=\linewidth]{test1_factor_ranking.png}
\caption{Factor ranking from representational probing. Scattered intensity
and crystallinity are the linearly decodable axes; orientation and peak count
are not.}
\label{fig:rank}
\end{figure}

\subsection{Causal ablations identify two drivers}

Table~\ref{tab:abl} and Figure~\ref{fig:abl} report re-inference ARI after
each input perturbation. A pass-through baseline reproduces the saved map
exactly ($\mathrm{ARI}=1.00$). Normalizing away the post-beam scattered
energy---removing \emph{how much} each position scatters while preserving the
relative spot pattern---collapses the map almost entirely
($\mathrm{ARI}=0.01$), establishing scattered intensity as a major driver.
Azimuthally averaging to the 1-D radial profile recovers only $\sim4$ coarse
groups ($\mathrm{ARI}=0.08$, $K\rightarrow4$): the full \emph{two}-
dimensional pattern is required for the remaining structure. Masking the
inner band degrades the map far more than masking the outer band
($\mathrm{ARI}=0.23$ vs.\ $0.53$), localizing the decision to \emph{low}~$q$;
blurring the sharp spots ($\sigma=2$) gives $\mathrm{ARI}=0.27$. We note one
artifact: per-frame min--max normalization leaves the map unchanged
($\mathrm{ARI}=1.00$) because the central beam pins each frame's maximum;
the post-beam ``scattered\_norm'' control above is the correct test.

\begin{table}[t]
\centering
\caption{Causal input ablations. ARI is computed against the original
$K=13$ map after re-inference. High ARI $=$ information not needed.}
\label{tab:abl}
\begin{tabular}{lcc}
\toprule
Ablation & ARI & $K$ \\
\midrule
Baseline (sanity)                 & 1.00 & 13 \\
Scattered-energy normalized       & \textbf{0.01} & 8 \\
Radial profile only (1-D)         & 0.08 & 4 \\
Low-$q$ masked                    & \textbf{0.23} & 11 \\
High-$q$ masked                   & 0.53 & 13 \\
Blur ($\sigma=2$)                 & 0.27 & 7 \\
Per-frame norm (artifact)         & 1.00 & 13 \\
\bottomrule
\end{tabular}
\end{table}

\begin{figure*}[t]
\centering
\includegraphics[width=0.92\textwidth]{test2_4_ablation_maps.png}
\caption{Class maps after each input ablation (re-inference). Removing the
overall scattered energy or reducing to the 1-D radial profile destroys the
partition; masking low~$q$ degrades it more than masking high~$q$.}
\label{fig:abl}
\end{figure*}

\subsection{Orientation is not a driver}

Indexing with a valid symmetry-derived (\texttt{auto}) library---rather than
a cubic-triangle plan that under-covers the triclinic IMC cell---confirms the
probing result. Across 367 confidently indexed positions spanning 43 distinct
(phase, zone-axis) labels, the DINO classes are essentially uncorrelated with
the recovered zone axis ($\mathrm{AMI}=0.05$, $\mathrm{ARI}=0.02$). The model
therefore does not key on absolute crystallographic orientation---the
expected outcome of the $\theta$-roll (azimuthal) augmentation, which trains
in-plane-rotation invariance.

\subsection{Spatial attribution corroborates low-$q$ dependence}

Grad-CAM and IG on the per-class average patterns
(Figure~\ref{fig:attr}) place attribution on the \emph{inner, low-$q$ rings
and spots}, with little weight in the high-$q$ corners---the spatial
counterpart of the low-$q$ ablation. The attended structure is
azimuthally distributed (it highlights rings rather than a single Bragg
direction), consistent with rotation invariance; more crystalline prototypes
show sharper low-$q$ attribution than diffuse, amorphous-like ones.

\begin{figure}[t]
\centering
\includegraphics[width=0.86\linewidth]{fig_paper_attribution.png}
\caption{Per-prototype class-average attribution. Columns: class mean,
Grad-CAM, integrated gradients, an example pattern, and its overlay.
Attention concentrates on the low-$q$ ring structure across prototypes.}
\label{fig:attr}
\end{figure}

\subsection{Per-class signatures}

The per-class mean diffraction patterns (Figure~\ref{fig:means}) form a
physically interpretable progression from diffuse, amorphous-like patterns to
sharp, ring-rich crystalline ones; per-class radial profiles separate mainly
by overall level and ring contrast, with finer two-dimensional differences
not visible in 1-D (explaining the collapse of the radial-only ablation to
$\sim4$ groups). Class-wise factor medians order primarily along the
scattered-intensity/crystallinity axis rather than ACOM correlation.

\begin{figure*}[t]
\centering
\includegraphics[width=0.92\textwidth]{class_mean_patterns.png}
\caption{Per-class mean NBED patterns ($K=13$). The classes span a
progression from diffuse/amorphous-like to sharp, crystalline ring patterns,
mirroring the scattered-intensity/crystallinity ordering.}
\label{fig:means}
\end{figure*}

\section{Conclusions}

The self-supervised partition of this IMC film rests on \emph{two} physical
axes, not one. The dominant, coarse axis is the total scattered/diffracted
intensity---a coupled crystallinity and mass-thickness signal whose removal
collapses the map (ARI $1.00\rightarrow0.01$). The fine $K=13$ structure is
supplied by the two-dimensional, low-$q$ diffraction pattern (rings/spots),
since the 1-D radial profile alone recovers only $\sim4$ groups and low-$q$
masking is far more destructive than high-$q$ masking. The partition is
\emph{not} based on crystallographic orientation (ACOM zone-axis
$\mathrm{AMI}=0.05$) or Bragg-peak count, and is invariant to in-plane
rotation---all consistent with the radial-gated, $\theta$-roll training
design. A practical caveat for downstream use: ``scattered intensity'' bundles
crystallinity with thickness/mass, which these measurements cannot separate;
an independent thickness map (e.g.\ a $t/\lambda$ estimate) would be required
to disentangle them.

\begin{acknowledgement}
Interpretability analysis (probing, ablation, ACOM cross-check, and
attribution) was carried out on the frozen, previously validated model
checkpoint.
\end{acknowledgement}

\begin{thebibliography}{9}

\bibitem{caron2021}
Caron, M.; Touvron, H.; Misra, I.; J\'egou, H.; Mairal, J.; Bojanowski, P.;
Joulin, A. Emerging Properties in Self-Supervised Vision Transformers.
\emph{Proc.\ IEEE/CVF Int.\ Conf.\ Comput.\ Vis.\ (ICCV)} \textbf{2021},
9650--9660.

\bibitem{savitzky2021}
Savitzky, B. H.\ et al. py4DSTEM: A Software Package for Four-Dimensional
Scanning Transmission Electron Microscopy Data Analysis. \emph{Microsc.\
Microanal.} \textbf{2021}, \emph{27}, 712--743.

\bibitem{ophus2022}
Ophus, C.\ et al. Automated Crystal Orientation Mapping in py4DSTEM Using
Sparse Correlation Matching. \emph{Microsc.\ Microanal.} \textbf{2022},
\emph{28}, 390--403.

\bibitem{khosla2020}
Khosla, P.; Teterwak, P.; Wang, C.; Sarna, A.; Tian, Y.; Isola, P.;
Maschinot, A.; Liu, C.; Krishnan, D. Supervised Contrastive Learning.
\emph{Adv.\ Neural Inf.\ Process.\ Syst.\ (NeurIPS)} \textbf{2020},
\emph{33}, 18661--18673.

\end{thebibliography}

\end{document}
